% ============================================================
%  Presentación Beamer "UNI Oscuro" — Métodos Numéricos (MB536)
%  Examen Final 2025-1 — Resuelto paso a paso
%  Compilar con: xelatex (dos pasadas)
% ============================================================
\documentclass[aspectratio=169,10pt]{beamer}

\usepackage{fontspec}
\usepackage[spanish,es-noshorthands]{babel}
\usepackage{tikz}
\usetikzlibrary{shapes.geometric,positioning}
\usepackage[most]{tcolorbox}
\usepackage{fontawesome5}
\usepackage{booktabs,colortbl,array,ragged2e,graphicx}
\usepackage{amsmath,amssymb}
\usepackage{mathtools}
\setlength{\emergencystretch}{3em}

% ---------- Paleta ----------
\definecolor{FondoOscuro}{HTML}{040812}
\definecolor{AzulPanel}{HTML}{0C111C}
\definecolor{Granate}{HTML}{660F1A}
\definecolor{GranateOscuro}{HTML}{3D0710}
\definecolor{GranateVelo}{HTML}{381520}
\definecolor{Teal}{HTML}{00A6A6}
\definecolor{TealOscuro}{HTML}{01565B}
\definecolor{Dorado}{HTML}{D4AF37}
\definecolor{Naranja}{HTML}{B46A35}
\definecolor{Cian}{HTML}{00F2FF}
\definecolor{Crema}{HTML}{FAF9F7}
\definecolor{CremaGris}{HTML}{F6F5F3}
\definecolor{TintaTexto}{HTML}{2F3E46}
\definecolor{TealSuave}{HTML}{F2FBFB}
\definecolor{GrisLinea}{HTML}{E2E1E0}
\definecolor{IconoTenue}{HTML}{0B2430}

% ---------- Fuentes ----------
\setsansfont[
  BoldFont=FiraSans-SemiBold.otf,
  ItalicFont=FiraSans-Italic.otf,
  BoldItalicFont=FiraSans-SemiBoldItalic.otf
]{FiraSans-Regular.otf}
\setmonofont[Scale=0.9]{FiraCode-Regular.ttf}
\usefonttheme{professionalfonts}
\renewcommand{\familydefault}{\sfdefault}

% ---------- METADATOS ----------
\newcommand{\sellocorto}{UNI | FIM | Métodos Numéricos}
\newcommand{\pietitulo}{Final 2025-1 · Resuelto}

% ---------- Ajustes globales beamer ----------
\setbeamertemplate{navigation symbols}{}
\setbeamersize{text margin left=8mm, text margin right=8mm}
\setbeamercolor{normal text}{fg=TintaTexto, bg=Crema}
\setbeamercolor{structure}{fg=Granate}
\setbeamercolor{alerted text}{fg=Granate}
\setbeamercolor{example text}{fg=TealOscuro}
\setbeamertemplate{itemize item}{\color{Teal}\raisebox{0.4ex}{\scriptsize\faCaretRight}}
\setbeamertemplate{itemize subitem}{\color{Granate}\raisebox{0.4ex}{\tiny\faCaretRight}}
\setbeamerfont{frametitle}{series=\bfseries,size=\large}

% ---------- Fondo de diapositivas de contenido ----------
\setbeamertemplate{background canvas}{%
\begin{tikzpicture}
  \fill[Crema] (0,0) rectangle (16,9);
  \draw[white,       line width=1.3pt] (0,2.1) .. controls (4.2,3.5) and (9.5,1.1) .. (16,3.1);
  \draw[GrisLinea!60,line width=0.9pt] (0,1.0) .. controls (5.0,2.3) and (10.5,0.3) .. (16,1.7);
  \draw[white,       line width=1.3pt] (0,5.3) .. controls (5.2,6.8) and (11.0,4.5) .. (16,6.3);
  \draw[GrisLinea!60,line width=0.9pt] (0,7.3) .. controls (6.0,8.5) and (11.5,6.5) .. (16,7.9);
\end{tikzpicture}}

% ---------- Cabecera ----------
\setbeamertemplate{frametitle}{%
\begin{tikzpicture}[remember picture, overlay]
  \fill[FondoOscuro] (current page.north west) rectangle ([yshift=-0.85cm]current page.north east);
  \fill[Granate] (current page.north west)
      -- ([xshift=7.6cm]current page.north west)
      -- ([xshift=6.7cm,yshift=-0.85cm]current page.north west)
      -- ([yshift=-0.85cm]current page.north west) -- cycle;
  \fill[Dorado] ([yshift=-0.91cm]current page.north west)
      rectangle ([xshift=6.7cm,yshift=-0.85cm]current page.north west);
  \fill[Teal] ([xshift=6.7cm,yshift=-0.97cm]current page.north west)
      rectangle ([yshift=-0.85cm]current page.north east);
  \fill[Teal] ([xshift=-0.42cm]current page.north east)
      rectangle ([xshift=-0.28cm,yshift=-0.30cm]current page.north east);
  \node[anchor=west, text=white] at ([xshift=0.55cm,yshift=-0.43cm]current page.north west)
      {\large\bfseries\insertframetitle};
\end{tikzpicture}%
\vspace*{0.42cm}}

% ---------- Pie de página ----------
\setbeamertemplate{footline}{%
\begin{tikzpicture}[remember picture, overlay]
  \fill[CremaGris] ([yshift=0.5cm]current page.south west) rectangle (current page.south east);
  \draw[Teal, line width=0.7pt] ([yshift=0.5cm]current page.south west) -- ([yshift=0.5cm]current page.south east);
  \fill[Granate] ([xshift=-0.34cm]current page.south east) rectangle ([yshift=0.5cm]current page.south east);
  \node[anchor=west, text=FondoOscuro] at ([xshift=0.55cm,yshift=0.25cm]current page.south west)
      {\scriptsize \faUniversity\hspace{1mm}\textbf{\sellocorto}\hspace{1.4mm}{\color{Teal}\faAngleRight}\hspace{1.4mm}\textit{\pietitulo}};
  \node[anchor=east, text=FondoOscuro] at ([xshift=-0.55cm,yshift=0.25cm]current page.south east)
      {\scriptsize\textbf{\insertframenumber}\hspace{0.8mm}{\tiny\inserttotalframenumber}};
\end{tikzpicture}}

% ---------- Tarjetas ----------
\newtcolorbox{cardidea}[3][]{enhanced, colback=white, frame hidden, boxrule=0pt,
  arc=1.6mm, left=3mm, right=3mm, top=2.2mm, bottom=2.4mm,
  borderline west={2.6pt}{0pt}{Granate},
  drop fuzzy shadow={black!30},
  fontupper=\small, coltext=TintaTexto,
  before upper={{\color{Granate}#3}\hspace{1.6mm}{\normalsize\bfseries\color{FondoOscuro}#2}\par\vspace{1.3mm}}}

\newtcolorbox{cardgear}[3][]{enhanced, colback=TealSuave, frame hidden, boxrule=0pt,
  arc=1.6mm, left=3mm, right=3mm, top=2.2mm, bottom=2.4mm,
  borderline west={2.6pt}{0pt}{Teal},
  drop fuzzy shadow={black!30},
  fontupper=\small, coltext=TintaTexto,
  before upper={{\color{TealOscuro}#3}\hspace{1.6mm}{\normalsize\bfseries\color{FondoOscuro}#2}\par\vspace{1.3mm}}}

\newtcolorbox{cardnota}[1][\faExclamationTriangle]{enhanced, colback=white,
  colframe=GrisLinea, boxrule=0.5pt, arc=1.4mm,
  left=3mm, right=3mm, top=1.8mm, bottom=2mm,
  drop fuzzy shadow={black!25},
  fontupper=\small, coltext=TintaTexto,
  before upper={{\color{FondoOscuro}#1}\hspace{1.6mm}}}

\newtcolorbox{cardlista}{enhanced, colback=white, frame hidden, boxrule=0pt,
  arc=1.2mm, left=3.5mm, right=2.5mm, top=2.4mm, bottom=2.4mm,
  borderline west={3pt}{0pt}{FondoOscuro},
  drop fuzzy shadow={black!30},
  fontupper=\small, coltext=Granate}

% ---------- Leyendas ----------
\newcommand{\tcap}[1]{\begin{center}\vspace{-1mm}{\scriptsize{\color{Granate}\bfseries Tabla.} {\color{TintaTexto!75}#1}}\end{center}\vspace{-2.5mm}}
\newcommand{\fcap}[1]{{\par\centering\scriptsize{\color{Granate}\bfseries Figura.} {\color{TintaTexto!75}#1}\par}}
\newcommand{\fsub}[1]{{\par\centering\scriptsize\color{TintaTexto!75}#1\par}}

% ---------- Portada de sección ----------
\newcommand{\portadaseccion}[4]{%
\begin{frame}[plain]
\begin{tikzpicture}[remember picture, overlay]
  \fill[FondoOscuro] (current page.south west) rectangle (current page.north east);
  \node[color=IconoTenue] at ([xshift=-3.4cm,yshift=0.3cm]current page.east)
      {\fontsize{62}{62}\selectfont #4};
  \fill[GranateVelo] (current page.north west)
      -- ([xshift=8.6cm]current page.north west)
      -- ([xshift=11.3cm]current page.south west)
      -- (current page.south west) -- cycle;
  \fill[GranateOscuro, opacity=0.75] (current page.north west)
      -- ([xshift=6.4cm]current page.north west)
      -- ([xshift=9.1cm]current page.south west)
      -- (current page.south west) -- cycle;
  \draw[Teal, line width=1.5pt] ([xshift=8.6cm]current page.north west)
      -- ([xshift=11.3cm]current page.south west);
  \draw[Naranja, line width=0.9pt] ([xshift=7.4cm]current page.north west)
      -- ([xshift=10.1cm]current page.south west);
  \node[anchor=west, text=white] at ([xshift=1.1cm,yshift=0.75cm]current page.west)
      {\fontsize{24}{28}\selectfont\bfseries #1.~#2};
  \draw[white, line width=0.9pt] ([xshift=1.15cm,yshift=0.12cm]current page.west) -- ++(4.8cm,0);
  \node[anchor=west, text=white] at ([xshift=1.15cm,yshift=-0.55cm]current page.west)
      {{\color{white}\faAngleDoubleRight}\hspace{1.5mm}\small #3};
\end{tikzpicture}
\end{frame}}

\begin{document}

% ============================================================
%  PORTADA
% ============================================================
\begin{frame}[plain]
\begin{tikzpicture}[remember picture, overlay]
  \fill[FondoOscuro] (current page.south west) rectangle (current page.north east);
  % --- fórmulas decorativas tenues ---
  \node[color=AzulPanel!80!white] at ([xshift=-4.4cm,yshift=-1.2cm]current page.north east)
      {\fontsize{17}{17}\selectfont $\displaystyle x^{(1)}=x^{(0)}-J_F^{-1}F$};
  \node[color=AzulPanel!80!white] at ([xshift=3.7cm,yshift=1.1cm]current page.south west)
      {\fontsize{19}{19}\selectfont $\det J=L_1L_2\sin\theta_2$};
  \node[color=AzulPanel!80!white] at ([xshift=-2.7cm,yshift=2.6cm]current page.south east)
      {\fontsize{18}{18}\selectfont $\tfrac{3h}{8}(f_0+3f_1+3f_2+f_3)$};
  % --- circuito decorativo ---
  \draw[Dorado, line width=0.8pt] ([xshift=-2.7cm,yshift=-3.4cm]current page.north east)
      -- ++(-1.5,-1.5) -- ++(-1.8,0);
  \fill[Cian] ([xshift=-6.0cm,yshift=-4.9cm]current page.north east) circle (0.055cm);
  \draw[Teal, line width=0.8pt] ([xshift=-1.6cm,yshift=-5.6cm]current page.north east)
      -- ++(-1.2,-1.2) -- ++(-2.2,0);
  \fill[Cian] ([xshift=-5.0cm,yshift=-6.8cm]current page.north east) circle (0.055cm);
  \fill[Cian] ([xshift=-1.6cm,yshift=-5.6cm]current page.north east) circle (0.05cm);
  % --- hexágono con ícono ---
  \node[regular polygon, regular polygon sides=6, minimum size=2.5cm,
        draw=Teal, line width=1.3pt, shape border rotate=30]
        at ([xshift=-3.1cm,yshift=-0.2cm]current page.east) {};
  \node[color=Teal] at ([xshift=-3.1cm,yshift=-0.2cm]current page.east)
      {\fontsize{26}{26}\selectfont \faClipboardCheck};
  % --- textos ---
  \node[anchor=north west, text=white] at ([xshift=1.0cm,yshift=-0.85cm]current page.north west)
      {\footnotesize\bfseries UNIVERSIDAD NACIONAL DE INGENIERIA\ \textbar\ FIM\ \textbar\ Métodos Numéricos (MB536)};
  \node[anchor=north west, text=white, align=left,
        font=\fontsize{19.5}{24}\selectfont\bfseries] at ([xshift=0.95cm,yshift=-1.55cm]current page.north west)
      {Examen Final 2025-1\\ \textemdash\ Resuelto};
  \node[anchor=north west, text=Teal] at ([xshift=1.0cm,yshift=-4.05cm]current page.north west)
      {\normalsize\itshape Sist.\ no lineales\ \textperiodcentered\ Cuadratura\ \textperiodcentered\ Interpolación\ \textperiodcentered\ Splines\ \textperiodcentered\ EDOs};
  \draw[Dorado, line width=0.6pt] ([xshift=1.0cm,yshift=2.45cm]current page.south west) -- ++(6.2cm,0);
  \node[anchor=south west, text=white] at ([xshift=1.0cm,yshift=1.75cm]current page.south west)
      {\normalsize\bfseries Problemas resueltos paso a paso};
  \node[anchor=south west, text=white] at ([xshift=1.0cm,yshift=1.25cm]current page.south west)
      {\small Fecha del examen: 11/07/25 \textbar\ Duración: 110 min \textbar\ 20 puntos};
  \node[anchor=south west, text=white!70!FondoOscuro] at ([xshift=1.0cm,yshift=0.78cm]current page.south west)
      {\footnotesize Coordinadora: Mag.\ Ing.\ Rosa Garrido Juárez \textperiodcentered\ Lima, 2025};
\end{tikzpicture}
\end{frame}

% ============================================================
%  MAPA
% ============================================================
\begin{frame}{Contenido del examen}
\begin{columns}[T]
\begin{column}{0.485\textwidth}
\begin{cardlista}
\textbf{Parte I} (preguntas cortas, 1.0 P c/u)\\[0.6mm]
P1.\; Sistema no lineal: Jacobiano y Newton\\[0.4mm]
P2.\; Completar script: trapecio y Simpson $1/3$\\[0.4mm]
P3.\; Diferencias divididas de Newton\\[0.4mm]
P4.\; Diferencias finitas en una EDO de frontera
\end{cardlista}
\end{column}
\begin{column}{0.485\textwidth}
\begin{cardlista}
\textbf{Parte II} (problemas largos, 4.0 P c/u)\\[0.6mm]
1.\; Punto fijo: contracción y 3 iteraciones\\[0.4mm]
2.\; Viscosidad: polinomio y spline cúbico\\[0.4mm]
3.\; Trabajo: Simpson $3/8$ vs.\ trapecio\\[0.4mm]
4.\; Polea: Euler modificado (Heun)
\end{cardlista}
\end{column}
\end{columns}
\vspace{2mm}
\begin{cardnota}[\faInfoCircle]
Se reproducen las \textbf{soluciones oficiales} del examen, \emph{expandidas}: se añade la aritmética intermedia, la justificación de cada paso y la verificación numérica. Los valores oficiales se conservan; las discrepancias detectadas al recalcular se \textbf{corrigen y se anotan}.
\end{cardnota}
\end{frame}

% ############################################################
%  PARTE I
% ############################################################
\portadaseccion{I}{Preguntas cortas}{Sistema no lineal, script de cuadratura, diferencias divididas y diferencias finitas}{\faListUl}

% ---- P1 enunciado ----
\begin{frame}{P1 — Enunciado (sistema no lineal)}
\begin{columns}[T]
\begin{column}{0.5\textwidth}
\begin{cardidea}{Enunciado \textmd{\small(1.0 P)}}{\faQuestionCircle}
{\footnotesize Dado el siguiente sistema de ecuaciones no lineales:
\[
f_1(x_1,x_2)=x_1^{3}+x_2-\tfrac12=0,
\]
\[
f_2(x_1,x_2)=x_1^{2}-x_2^{2}=0.
\]
\textbf{Complete la tabla} con: (i) el número de soluciones reales; (ii) el Jacobiano $J_F$ de la función vectorial $F=[f_1;f_2]$; (iii) la matriz inversa $J_F^{-1}$ evaluada en $[\tfrac12;\tfrac12]$; y (iv) la primera aproximación por el método de Newton con punto inicial $[\tfrac12;\tfrac12]$.}
\end{cardidea}
\end{column}
\begin{column}{0.48\textwidth}
\begin{cardnota}[\faInfoCircle]
{\footnotesize \textbf{Newton para sistemas:}
\[
x^{(1)}=x^{(0)}-J_F^{-1}(x^{(0)})\,F(x^{(0)}).
\]
Se requieren $J_F$, su inversa en el punto inicial y el vector $F(x^{(0)})$.}
\end{cardnota}
\vspace{1mm}
\begin{cardnota}[\faTable]
{\footnotesize \textbf{Casillas a completar:} $N^{\!o}$ soluciones reales $\;\vert\;$ $J_F(x_1,x_2)$ $\;\vert\;$ $J_F^{-1}(\tfrac12,\tfrac12)$ $\;\vert\;$ $x^{(1)}$.}
\end{cardnota}
\end{column}
\end{columns}
\end{frame}

% ---- P1 solución: Jacobiano e inversa ----
\begin{frame}{P1 — Jacobiano, determinante e inversa}
\begin{columns}[T]
\begin{column}{0.5\textwidth}
\begin{cardgear}{Jacobiano de $F=[f_1;f_2]$}{\faCogs}
{\footnotesize Derivando cada $f_i$ respecto de $x_1,x_2$:
\[
J_F(x_1,x_2)=\begin{pmatrix} 3x_1^{2} & 1\\[1mm] 2x_1 & -2x_2 \end{pmatrix}.
\]
Evaluado en $(\tfrac12,\tfrac12)$ ($3x_1^2=0.75$):
\[
J_F(\tfrac12,\tfrac12)=\begin{pmatrix} 0.75 & 1\\[1mm] 1 & -1 \end{pmatrix}.
\]
Determinante:
\[
\det J_F=(0.75)(-1)-(1)(1)=-\tfrac74=-1.75.
\]}
\end{cardgear}
\end{column}
\begin{column}{0.48\textwidth}
\begin{cardgear}{Matriz inversa en $(\tfrac12,\tfrac12)$}{\faSuperscript}
{\footnotesize Para $2\times2$: $A^{-1}=\tfrac1{\det}\begin{psmallmatrix} d & -b\\ -c & a\end{psmallmatrix}$.
\[
J_F^{-1}=\frac1{-7/4}\begin{pmatrix} -1 & -1\\ -1 & 0.75\end{pmatrix}
\]
\[
=\frac17\begin{pmatrix} 4 & 4\\ 4 & -3\end{pmatrix}.
\]}
\end{cardgear}
\vspace{1mm}
\begin{cardnota}[\faCheckCircle]
{\footnotesize $\tfrac{-1}{-7/4}=\tfrac47$ y $\tfrac{0.75}{-7/4}=-\tfrac37$: se factoriza $\tfrac17$ dejando enteros $\begin{psmallmatrix}4&4\\4&-3\end{psmallmatrix}$.}
\end{cardnota}
\end{column}
\end{columns}
\end{frame}

% ---- P1 solución: Newton + tabla ----
\begin{frame}{P1 — Número de soluciones y primera iteración de Newton}
\begin{columns}[T]
\begin{column}{0.48\textwidth}
\begin{cardgear}{Soluciones reales y evaluación de $F$}{\faCalculator}
{\footnotesize \textbf{$N^{\!o}$ soluciones reales: 2.}\ De $f_2=0$: $x_2=\pm x_1$; sustituyendo en $f_1$ se obtienen dos raíces reales.
\smallskip

$F$ en el punto inicial ($x_1^3=0.125$):
\[
F(\tfrac12,\tfrac12)=\begin{pmatrix}0.5^3+0.5-0.5\\ 0.5^2-0.5^2\end{pmatrix}=\begin{pmatrix}0.125\\ 0\end{pmatrix}.
\]}
\end{cardgear}
\end{column}
\begin{column}{0.5\textwidth}
\begin{cardgear}{Paso de Newton}{\faArrowRight}
{\footnotesize
\[
x^{(1)}=\begin{pmatrix}\tfrac12\\\tfrac12\end{pmatrix}-\frac17\begin{pmatrix}4&4\\4&-3\end{pmatrix}\begin{pmatrix}0.125\\0\end{pmatrix}.
\]
Producto: $\tfrac17(4\cdot0.125,\,4\cdot0.125)^{\!\top}=\tfrac17(0.5,0.5)^{\!\top}$.
\[
x^{(1)}=\begin{pmatrix}0.5\\0.5\end{pmatrix}-\begin{pmatrix}0.0714\\0.0714\end{pmatrix}=\tfrac37\begin{pmatrix}1\\1\end{pmatrix}.
\]}
\end{cardgear}
\vspace{1mm}
\begin{cardnota}[\faCheckCircle]
{\footnotesize \textbf{Resultado:}\ $x^{(1)}=\tfrac37(1,1)^{\!\top}\approx(0.4285,\,0.4285)^{\!\top}$.}
\end{cardnota}
\end{column}
\end{columns}
\end{frame}

% ---- P2 enunciado ----
\begin{frame}{P2 — Enunciado (energía por cuadratura)}
\begin{columns}[T]
\begin{column}{0.56\textwidth}
\begin{cardidea}{Enunciado \textmd{\small(1.0 P)}}{\faQuestionCircle}
{\footnotesize Se requiere calcular la \textbf{energía acumulada} en el intervalo $[0,\pi/2]$ de un proceso físico modelado por $f(x)=e^{x}\cos x$. La energía se obtiene integrando $f(x)$ en dicho intervalo:
\[
E=\int_0^{\pi/2} e^{x}\cos x\,dx.
\]
Implemente en \textbf{MATLAB} el método del \textbf{trapecio compuesto} y el método de \textbf{Simpson $1/3$ compuesto}, ambos con $n=4$ subintervalos. \textbf{Complete los espacios en blanco} del script (resaltados en granate).}
\end{cardidea}
\end{column}
\begin{column}{0.42\textwidth}
\begin{cardnota}[\faInfoCircle]
{\footnotesize Con $n=4$ y $[a,b]=[0,\pi/2]$ hay $5$ nodos $x_0,\dots,x_4$ y $h=\dfrac{b-a}{n}=\dfrac{\pi}{8}$.}
\end{cardnota}
\vspace{1mm}
\begin{cardnota}[\faKey]
{\footnotesize \textbf{Valor exacto} (por partes):
\[
E=\tfrac12\big(e^{\pi/2}-1\big)\approx 1.9052.
\]
Sirve de referencia para los errores.}
\end{cardnota}
\end{column}
\end{columns}
\end{frame}

% ---- P2 script ----
\begin{frame}{P2 — Script con los 6 blancos completados}
\begin{columns}[T]
\begin{column}{0.62\textwidth}
\begin{cardgear}{Script MATLAB (blancos en granate)}{\faCode}
{\scriptsize\ttfamily
clear all\\
\% Define anonymous function\\
f = @({\color{Granate}x}) {\color{Granate}exp(x).*cos(x)};\ \ \% 1\\[0.4mm]
\% Symbolic exact integral\\
syms x\\
I\_sym = int(f(x), {\color{Granate}0, pi/2});\ \ \% 2\\
I\_exact = double(I\_sym);\\
\% Composite Trapezoidal Rule\\
a = 0;\ b = pi/2;\ n = 4;\\
h = (b - a) / {\color{Granate}n};\ \ \% 3\\
xvals = [a:h:b];\\
I\_trap = (h/2)*(f(a) + 2*sum(f(xvals({\color{Granate}2:end-1}))) + f(b));\ \ \% 4\\[0.4mm]
\% Composite Simpson's Rule\\
xodd = xvals(2:2:n);\\
xeven = xvals({\color{Granate}3:2:n-1});\ \ \% 5\\
I\_simp = (h/3)*(f(a) + 4*sum(f(xodd)) + 2*sum(f(xeven)) + f(b));\\
\% Error comparison\\
error\_trap = abs(I\_trap - I\_exact);\\
error\_simp = abs({\color{Granate}I\_simp - I\_exact});\ \ \% 6
}
\end{cardgear}
\end{column}
\begin{column}{0.36\textwidth}
\begin{cardnota}[\faListOl]
{\scriptsize
\textbf{1.} \texttt{@(x) exp(x).*cos(x)} — función anónima; \texttt{.*} para producto elemento a elemento.\\[0.6mm]
\textbf{2.} \texttt{0, pi/2} — límites de la integral simbólica.\\[0.6mm]
\textbf{3.} \texttt{n} — paso $h=(b-a)/n$.\\[0.6mm]
\textbf{4.} \texttt{2:end-1} — nodos \emph{interiores} (peso $2$) en el trapecio.\\[0.6mm]
\textbf{5.} \texttt{3:2:n-1} — nodos de índice par (peso $2$) en Simpson.\\[0.6mm]
\textbf{6.} \texttt{I\_simp - I\_exact} — error de Simpson frente al exacto.}
\end{cardnota}
\end{column}
\end{columns}
\end{frame}

% ---- P2 verificación ----
\begin{frame}{P2 — Verificación numérica de las dos reglas}
\begin{columns}[T]
\begin{column}{0.52\textwidth}
\begin{cardgear}{Nodos y valores ($h=\pi/8$)}{\faTable}
\centering
{\footnotesize\renewcommand{\arraystretch}{1.2}\setlength{\tabcolsep}{6pt}\arrayrulecolor{Granate}
\begin{tabular}{c|c|c}
$i$ & $x_i$ & $f(x_i)=e^{x_i}\cos x_i$\\\hline
0 & 0.0000 & 1.0000\\
1 & 0.3927 & 1.3682\\
2 & 0.7854 & 1.5510\\
3 & 1.1781 & 1.2430\\
4 & 1.5708 & 0.0000\\
\end{tabular}}
\end{cardgear}
\end{column}
\begin{column}{0.46\textwidth}
\begin{cardgear}{Resultados}{\faCalculator}
{\scriptsize
\[
I_{\text{trap}}=\tfrac{h}{2}\!\left[f_0+2(f_1{+}f_2{+}f_3)+f_4\right]\approx1.8308.
\]
\[
I_{\text{simp}}=\tfrac{h}{3}\!\left[f_0+4(f_1{+}f_3)+2f_2+f_4\right]\approx1.9042.
\]}
\end{cardgear}
\vspace{1mm}
\begin{cardnota}[\faCheckCircle]
{\footnotesize $I_{\text{exact}}\approx1.9052$. Errores: $|{\cdot}|_{\text{trap}}\approx0.074$ y $|{\cdot}|_{\text{simp}}\approx0.001$: Simpson $1/3$ es mucho más preciso a igual $n$.}
\end{cardnota}
\end{column}
\end{columns}
\end{frame}

% ---- P3 enunciado + solución ----
\begin{frame}{P3 — Diferencias divididas de Newton}
\begin{columns}[T]
\begin{column}{0.42\textwidth}
\begin{cardidea}{Enunciado \textmd{\small(1.0 P)}}{\faQuestionCircle}
{\footnotesize Dado el conjunto de datos $x=[1.5,\ 2.0,\ 2.8]$ y $f(x)=[3.2,\ 4.5,\ 6.8]$, elabore la \textbf{tabla de diferencias divididas} hasta el segundo orden, obtenga el \textbf{polinomio interpolante de Newton de grado dos} y utilícelo para estimar el valor de $\mathbf{f(2.3)}$.}
\end{cardidea}
\vspace{1mm}
\begin{cardgear}{Tabla de diferencias divididas}{\faTable}
\centering
{\footnotesize\renewcommand{\arraystretch}{1.25}\setlength{\tabcolsep}{5pt}\arrayrulecolor{Granate}
\begin{tabular}{c|c|c|c}
$x_i$ & $f_i$ & $1^{er}$ & $2^{do}$\\\hline
1.5 & 3.2 & 2.6 & 0.2115\\
2.0 & 4.5 & 2.875 & \\
2.8 & 6.8 & & \\
\end{tabular}}
\end{cardgear}
\end{column}
\begin{column}{0.56\textwidth}
\begin{cardgear}{Cálculo de las diferencias}{\faCalculator}
{\footnotesize
\[
f[x_0,x_1]=\frac{4.5-3.2}{2.0-1.5}=\frac{1.3}{0.5}=2.6,
\]
\[
f[x_1,x_2]=\frac{6.8-4.5}{2.8-2.0}=\frac{2.3}{0.8}=2.875,
\]
\[
f[x_0,x_1,x_2]=\frac{2.875-2.6}{2.8-1.5}=\frac{0.275}{1.3}=0.2115.
\]}
\end{cardgear}
\vspace{1mm}
\begin{cardgear}{Polinomio y evaluación en $x=2.3$}{\faSuperscript}
{\footnotesize
\[
P_2(x)=3.2+2.6(x{-}1.5)+0.2115(x{-}1.5)(x{-}2.0).
\]
\[
P_2(2.3)=3.2+2.6(0.8)+0.2115(0.8)(0.3)
\]
\[
=3.2+2.08+0.05076=\boxed{5.331}.
\]}
\end{cardgear}
\end{column}
\end{columns}
\end{frame}

% ---- P4 enunciado ----
\begin{frame}{P4 — Enunciado (diferencias finitas en EDO de frontera)}
\begin{columns}[T]
\begin{column}{0.52\textwidth}
\begin{cardidea}{Enunciado \textmd{\small(1.0 P)}}{\faQuestionCircle}
{\footnotesize Considere la ecuación diferencial ordinaria con condiciones de frontera:
\[
y''+2y'+3y+4x=0,
\]
\[
y'(0)=1,\qquad y'(1)=2.
\]
Aplicando el método de \textbf{diferencias finitas} se obtiene el sistema lineal
\[
\begin{pmatrix}-5&a&b\\ c&-5&d\\ e&f&-5\end{pmatrix}\!\!\begin{pmatrix}y_0\\y_1\\y_2\end{pmatrix}=\begin{pmatrix}u\\v\\w\end{pmatrix}.
\]
\textbf{Complete los coeficientes faltantes} justificando su procedimiento.}
\end{cardidea}
\end{column}
\begin{column}{0.46\textwidth}
\begin{cardnota}[\faInfoCircle]
{\footnotesize \textbf{Malla ($h=0.5$).} Nodos internos $x_0{=}0,\ x_1{=}0.5,\ x_2{=}1$; nodos \textbf{ficticios} $y_{-1}$ y $y_3$ en los extremos para imponer las derivadas de frontera.}
\end{cardnota}
\vspace{1mm}
\begin{cardnota}[\faRuler]
{\footnotesize \textbf{Aproximaciones centradas:}
\[
y''\!\approx\!\tfrac{y_{i+1}-2y_i+y_{i-1}}{h^2},\quad
y'\!\approx\!\tfrac{y_{i+1}-y_{i-1}}{2h}.
\]}
\end{cardnota}
\end{column}
\end{columns}
\end{frame}

% ---- P4 solución 1: discretización ----
\begin{frame}{P4 — Discretización y coeficientes por nodo}
\begin{columns}[T]
\begin{column}{0.53\textwidth}
\begin{cardgear}{Fórmula discreta ($i=0,1,2$)}{\faCogs}
{\footnotesize Sustituyendo las diferencias centradas en $y''+2y'+3y=-4x$:
\[
\tfrac{y_{i+1}-2y_i+y_{i-1}}{h^2}+2\,\tfrac{y_{i+1}-y_{i-1}}{2h}+3y_i=-4x_i.
\]
Con $h=0.5$: $\tfrac1{h^2}=4$, $\tfrac1{h}=2$. Agrupando por nodo:
\[
\underbrace{(4-2)}_{2}y_{i-1}+\underbrace{(-8+3)}_{-5}y_i+\underbrace{(4+2)}_{6}y_{i+1}=-4x_i.
\]}
\end{cardgear}
\vspace{1mm}
\begin{cardnota}[\faCheckCircle]
{\footnotesize Esquema base: $\ 2y_{i-1}-5y_i+6y_{i+1}=-4x_i.$}
\end{cardnota}
\end{column}
\begin{column}{0.45\textwidth}
\begin{cardgear}{Condiciones de frontera}{\faTable}
{\footnotesize Con derivadas centradas ($2h=1$):
\[
\frac{y_1-y_{-1}}{2h}=1\ \Rightarrow\ y_{-1}=y_1-1,
\]
\[
\frac{y_3-y_1}{2h}=2\ \Rightarrow\ y_3=y_1+2.
\]
Estas relaciones \emph{eliminan} los nodos ficticios $y_{-1}$ y $y_3$.}
\end{cardgear}
\end{column}
\end{columns}
\end{frame}

% ---- P4 solución 2: sistema ----
\begin{frame}{P4 — Ensamblaje del sistema lineal}
\begin{columns}[T]
\begin{column}{0.55\textwidth}
\begin{cardgear}{Ecuación en cada nodo}{\faCalculator}
{\footnotesize \textbf{$i=0$} ($x_0=0$): $2y_{-1}-5y_0+6y_1=0$. Con $y_{-1}=y_1-1$:
\[
-5y_0+8y_1=2.
\]
\textbf{$i=1$} ($x_1=0.5$): $2y_0-5y_1+6y_2=-2.$
\smallskip

\textbf{$i=2$} ($x_2=1$): $2y_1-5y_2+6y_3=-4$. Con $y_3=y_1+2$:
\[
8y_1-5y_2=-16.
\]}
\end{cardgear}
\end{column}
\begin{column}{0.43\textwidth}
\begin{cardidea}{Sistema final}{\faSuperscript}
{\footnotesize
\[
\begin{pmatrix}-5&8&0\\ 2&-5&6\\ 0&8&-5\end{pmatrix}\!\!\begin{pmatrix}y_0\\y_1\\y_2\end{pmatrix}=\begin{pmatrix}2\\-2\\-16\end{pmatrix}.
\]}
\end{cardidea}
\vspace{1mm}
\begin{cardnota}[\faCheckCircle]
{\footnotesize Coeficientes: $a=8,\ b=0,\ c=2,\ d=6,\ e=0,\ f=8$ y RHS $u=2,\ v=-2,\ w=-16$.}
\end{cardnota}
\end{column}
\end{columns}
\end{frame}

% ############################################################
%  PARTE II — PROBLEMA 1
% ############################################################
\portadaseccion{II}{Problema 1 — Aproximaciones sucesivas}{Esquema de punto fijo, criterio de contracción y 3 iteraciones}{\faChartLine}

% ---- PII-1 enunciado ----
\begin{frame}{Problema 1 — Enunciado (punto fijo)}
\begin{columns}[T]
\begin{column}{0.5\textwidth}
\begin{cardidea}{Sistema no lineal}{\faCogs}
{\footnotesize Sea el siguiente sistema de ecuaciones no lineales:
\[
f_1(x,y)=x^{2}-20x+y^{2}+8,
\]
\[
f_2(x,y)=xy^{2}+x-20y+8.
\]
El \textbf{gráfico} del enunciado muestra las dos curvas $f_1=0$ y $f_2=0$ en $[0,2]\times[0,2]$; se cruzan cerca de $(0.42,\,0.42)$, que es el \textbf{valor inicial} sugerido $(x_0,y_0)=(0.5,\,0.5)$.}
\end{cardidea}
\end{column}
\begin{column}{0.48\textwidth}
\begin{cardidea}{Se pide}{\faTasks}
{\footnotesize
\textbf{a)\ (1.0 P)}\ Plantee un esquema iterativo de \textbf{aproximaciones sucesivas}, garantizando el cumplimiento del criterio de convergencia. Comente acerca de la tasa de convergencia.\\[1.0mm]
\textbf{b)\ (2.0 P)}\ Realice \textbf{03 iteraciones} del método de aproximaciones sucesivas a partir del valor inicial mostrado en el gráfico.\\[1.0mm]
\textbf{c)\ (1.0 P)}\ Calcule el \textbf{error relativo porcentual} en cada iteración usando la \textbf{norma infinita}. Analice si los resultados coinciden con lo concluido en el ítem a).}
\end{cardidea}
\end{column}
\end{columns}
\end{frame}

% ---- PII-1(a) ----
\begin{frame}{Problema 1(a) — Esquema y criterio de contracción}
\begin{columns}[T]
\begin{column}{0.5\textwidth}
\begin{cardgear}{Despeje del esquema $x=G(x,y)$}{\faSuperscript}
{\footnotesize De $f_1=0$ se despeja $20x$; de $f_2=0$ se despeja $20y$:
\[
x=G_1(x,y)=\frac{x^{2}+y^{2}+8}{20},
\]
\[
y=G_2(x,y)=\frac{xy^{2}+x+8}{20}.
\]
Jacobiano del iterador:
\[
J=\begin{pmatrix}\dfrac{x}{10} & \dfrac{y}{10}\\[2mm] \dfrac{y^{2}}{20}+\dfrac1{20} & \dfrac{xy}{10}\end{pmatrix}.
\]}
\end{cardgear}
\end{column}
\begin{column}{0.48\textwidth}
\begin{cardgear}{Contracción en $(0.5,0.5)$}{\faRulerHorizontal}
{\footnotesize
\[
J(0.5,0.5)=\begin{pmatrix}0.05 & 0.05\\ 0.0625 & 0.025\end{pmatrix}.
\]
Sumas por fila: $0.05+0.05=0.10$ y $0.0625+0.025=0.0875$.
\[
\lVert J\rVert_\infty=\max(0.10,\,0.0875)=0.1<1.
\]}
\end{cardgear}
\vspace{1mm}
\begin{cardnota}[\faCheckCircle]
{\footnotesize El arreglo es \textbf{convergente} por ser $\lVert J\rVert_\infty<1$. Al estar muy cerca de $0$, la \textbf{convergencia es rápida}.}
\end{cardnota}
\end{column}
\end{columns}
\end{frame}

% ---- PII-1(b)(c) ----
\begin{frame}{Problema 1(b, c) — Iteraciones y error relativo}
\begin{columns}[T]
\begin{column}{0.54\textwidth}
\begin{cardgear}{Tres iteraciones desde $(0.5,0.5)$}{\faTable}
{\footnotesize Iterando $x_{n+1}=\dfrac{x_n^2+y_n^2+8}{20}$,\ $y_{n+1}=\dfrac{x_ny_n^2+x_n+8}{20}$:}\\[1mm]
\centering
{\footnotesize\renewcommand{\arraystretch}{1.25}\setlength{\tabcolsep}{5pt}\arrayrulecolor{Granate}
\begin{tabular}{c|c|c|c}
$n$ & $x_n$ & $y_n$ & $E_{\text{rel}}\,(\%)$\\\hline
0 & 0.50000 & 0.50000 & --\\
1 & 0.42500 & 0.43125 & 17.39\\
2 & 0.41833 & 0.42520 & 1.57\\
3 & 0.41779 & 0.42470 & 0.13\\
\end{tabular}}
\end{cardgear}
\end{column}
\begin{column}{0.44\textwidth}
\begin{cardgear}{Error con norma $\infty$}{\faCalculator}
{\footnotesize
\[
E_{\text{rel}}=\frac{\lVert x_{n+1}-x_n\rVert_\infty}{\lVert x_{n+1}\rVert_\infty}\times100\%.
\]
Ej.\ $n{=}1$: $\tfrac{|{-}0.075|}{0.43125}=17.39\%$.}
\end{cardgear}
\vspace{1mm}
\begin{cardnota}[\faChartLine]
{\footnotesize El error \textbf{decrece} $17.39\to1.57\to0.13\%$: el método es \textbf{convergente} y \emph{rápido}, coherente con $\lVert J\rVert_\infty=0.1$ del ítem (a).}
\end{cardnota}
\end{column}
\end{columns}
\end{frame}

% ############################################################
%  PARTE II — PROBLEMA 2
% ############################################################
\portadaseccion{II}{Problema 2 — Viscosidad}{Polinomio interpolante cúbico y spline cúbico natural con datos ruidosos}{\faWater}

% ---- PII-2 enunciado ----
\begin{frame}{Problema 2 — Enunciado (viscosidad de un lubricante)}
\begin{columns}[T]
\begin{column}{0.5\textwidth}
\begin{cardidea}{Contexto}{\faInfoCircle}
{\footnotesize En un laboratorio de ingeniería de materiales se analiza cómo varía la \textbf{viscosidad cinemática} $\mu_k$ (en $\mathrm{m^2/s}\times10^{-6}$) de un lubricante en función de la temperatura $T$ ($^{\circ}$C). Los datos experimentales, con cierto \textbf{ruido} inherente, son:}\\[0.8mm]
\centering
{\footnotesize\renewcommand{\arraystretch}{1.25}\setlength{\tabcolsep}{7pt}\arrayrulecolor{Granate}
\begin{tabular}{c|cccc}
$T$ ($^{\circ}$C) & 0 & 20 & 40 & 50\\\hline
$\mu_k$ & 1.787 & 1.003 & 0.659 & 0.547\\
\end{tabular}}
\end{cardidea}
\end{column}
\begin{column}{0.48\textwidth}
\begin{cardidea}{Se pide}{\faTasks}
{\footnotesize
\textbf{a)\ (1.0 P)}\ De la tabla se obtuvo el polinomio interpolante cúbico $P(T)=a(T{-}20)(T{-}40)(T{-}50)+b\,T(T{-}40)(T{-}50)+c\,T(T{-}20)(T{-}50)+d\,T(T{-}20)(T{-}40)$. Halle $a,b,c,d$.\\[1.0mm]
\textbf{b)\ (2.0 P)}\ Calcule el \textbf{spline cúbico natural} para estimar la viscosidad cinemática a $35\,^{\circ}$C.\\[1.0mm]
\textbf{c)\ (1.0 P)}\ Con el spline del inciso (b), calcule la \textbf{razón de cambio instantánea} de $\mu_k$ respecto a $T$ cuando $T=20\,^{\circ}$C. Discuta las ventajas de splines cúbicos frente a un único polinomio de alto grado.}
\end{cardidea}
\end{column}
\end{columns}
\end{frame}

% ---- PII-2(a) ----
\begin{frame}{Problema 2(a) — Polinomio interpolante cúbico}
\begin{columns}[T]
\begin{column}{0.52\textwidth}
\begin{cardgear}{Forma producto de Lagrange}{\faSuperscript}
{\scriptsize Cada coeficiente se aísla evaluando $P$ en su propio nodo (los demás términos se anulan): $a=\dfrac{P(0)}{(0{-}20)(0{-}40)(0{-}50)}$, etc.
\[
a=\frac{1.787}{(-20)(-40)(-50)}=\frac{1.787}{-40000}=-4.4675\!\times\!10^{-5},
\]
\[
b=\frac{1.003}{(20)(-20)(-30)}=\frac{1.003}{12000}=8.3583\!\times\!10^{-5}.
\]}
\end{cardgear}
\end{column}
\begin{column}{0.46\textwidth}
\begin{cardgear}{Restantes}{\faCalculator}
{\scriptsize
\[
c=\frac{0.659}{(40)(20)(-10)}=\frac{0.659}{-8000}=-8.2375\!\times\!10^{-5},
\]
\[
d=\frac{0.547}{(50)(30)(10)}=\frac{0.547}{15000}=3.6467\!\times\!10^{-5}.
\]}
\end{cardgear}
\vspace{1mm}
\begin{cardnota}[\faCheckCircle]
{\scriptsize \textbf{Resultado:}\ $a=-4.4675\!\times\!10^{-5}$,\ $b=8.3583\!\times\!10^{-5}$,\ $c=-8.2375\!\times\!10^{-5}$,\ $d=3.6467\!\times\!10^{-5}$.}
\end{cardnota}
\end{column}
\end{columns}
\end{frame}

% ---- PII-2(b) sistema tridiagonal ----
\begin{frame}{Problema 2(b) — Spline natural: sistema de momentos}
\begin{columns}[T]
\begin{column}{0.5\textwidth}
\begin{cardgear}{Pasos y pendientes}{\faTable}
{\footnotesize $h_1=h_2=20$, $h_3=10$. Diferencias divididas de $1^{er}$ orden:
\[
d_1=\tfrac{1.003-1.787}{20}=-0.0392,
\]
\[
d_2=\tfrac{0.659-1.003}{20}=-0.0172,
\]
\[
d_3=\tfrac{0.547-0.659}{10}=-0.0112.
\]
Spline \emph{natural}: $M_0=M_3=0$.}
\end{cardgear}
\end{column}
\begin{column}{0.48\textwidth}
\begin{cardgear}{Sistema $2\times2$ en $M_1,M_2$}{\faCalculator}
{\scriptsize Ecuación de momentos $h_{i}M_{i-1}+2(h_i{+}h_{i+1})M_i+h_{i+1}M_{i+1}=6(d_{i+1}{-}d_i)$:
\[
\begin{pmatrix}80&20\\ 20&60\end{pmatrix}\!\!\begin{pmatrix}M_1\\M_2\end{pmatrix}=\begin{pmatrix}0.132\\ 0.036\end{pmatrix}.
\]
Solución:
\[
M_1=1.6364\!\times\!10^{-3},\ \ M_2=5.4545\!\times\!10^{-5}.
\]}
\end{cardgear}
\vspace{1mm}
\begin{cardnota}[\faExclamationTriangle]
{\scriptsize \textbf{Corrección.} La clave imprime la $2^{\text{a}}$ componente como $0.0366$ (y $M_2\!\approx\!6.5\!\times\!10^{-5}$); al recalcular $6(d_3{-}d_2)=6(0.006)=0.036$, con lo que $M_2=5.45\!\times\!10^{-5}$ — valor coherente con los coeficientes y con $\mu(35)$.}
\end{cardnota}
\end{column}
\end{columns}
\end{frame}

% ---- PII-2(b) coeficientes y evaluación ----
\begin{frame}{Problema 2(b) — Coeficientes por tramo y $\mu(35)$}
\begin{columns}[T]
\begin{column}{0.56\textwidth}
\begin{cardgear}{Coeficientes $S_i(T)=a_i\Delta^3+b_i\Delta^2+c_i\Delta+d_i$}{\faTable}
{\scriptsize $a_i=\tfrac{M_{i+1}-M_i}{6h_i}$, $b_i=\tfrac{M_i}{2}$, $c_i=d_i^{\text{pend}}-\tfrac{h_i(2M_i+M_{i+1})}{6}$, $d_i=\mu_i$:}\\[1mm]
\centering
{\scriptsize\renewcommand{\arraystretch}{1.3}\setlength{\tabcolsep}{4pt}\arrayrulecolor{Granate}
\begin{tabular}{c|cccc}
Tramo & $a_i$ & $b_i$ & $c_i$ & $d_i$\\\hline
$[0,20]$ & $1.3636\text{e-}5$ & $0$ & $-0.044655$ & $1.787$\\
$[20,40]$ & $-1.3182\text{e-}5$ & $8.182\text{e-}4$ & $-0.028291$ & $1.003$\\
$[40,50]$ & $-9.091\text{e-}7$ & $2.727\text{e-}5$ & $-0.011382$ & $0.659$\\
\end{tabular}}
\end{cardgear}
\end{column}
\begin{column}{0.42\textwidth}
\begin{cardgear}{Evaluación en $T=35$}{\faCalculator}
{\footnotesize $T=35\in[20,40]$, $\Delta=T-20=15$:
\[
S_1(35)=a_1(15)^3+b_1(15)^2+c_1(15)+d_1.
\]
\[
=-0.0445+0.1841-0.4244+1.003.
\]}
\end{cardgear}
\vspace{1mm}
\begin{cardnota}[\faCheckCircle]
{\footnotesize \textbf{Resultado:}\ $\mu(35)\approx0.718\ (\times10^{-6}\ \mathrm{m^2/s})$.}
\end{cardnota}
\end{column}
\end{columns}
\end{frame}

% ---- PII-2(c) ----
\begin{frame}{Problema 2(c) — Razón de cambio y ventajas del spline}
\begin{columns}[T]
\begin{column}{0.48\textwidth}
\begin{cardgear}{Derivada en $T=20$}{\faSuperscript}
{\footnotesize Sobre el tramo $[0,20]$, $S_0(T)=a_0(T)^3+c_0T+1.787$:
\[
S_0'(T)=3a_0T^{2}+c_0.
\]
En $T=20$ ($a_0=1.3636\!\times\!10^{-5}$, $c_0=-0.044655$):
\[
S_0'(20)=3(1.3636\text{e-}5)(400)-0.044655
\]
\[
\approx-2.83\!\times\!10^{-8}\ \tfrac{\mathrm{m^2/s}}{^{\circ}\!C}\ (\times10^{-6}).
\]}
\end{cardgear}
\end{column}
\begin{column}{0.5\textwidth}
\begin{cardidea}{Ventajas de los splines cúbicos}{\faLightbulb}
{\footnotesize
\begin{itemize}
\item \textbf{Evitan oscilaciones} de Runge, típicas de un único polinomio de alto grado.
\item Mayor \textbf{estabilidad numérica} y precisión \emph{local} tramo a tramo.
\item \textbf{Continuidad $C^2$}: primera y segunda derivadas continuas $\Rightarrow$ curva suave.
\item Muy adecuados para \textbf{datos ruidosos} como los del laboratorio.
\end{itemize}}
\end{cardidea}
\end{column}
\end{columns}
\end{frame}

% ############################################################
%  PARTE II — PROBLEMA 3
% ############################################################
\portadaseccion{II}{Problema 3 — Trabajo por fricción}{Integración numérica: Simpson $3/8$ frente a la regla del trapecio}{\faWater}

% ---- PII-3 enunciado ----
\begin{frame}{Problema 3 — Enunciado (trabajo de la fuerza de fricción)}
\begin{columns}[T]
\begin{column}{0.52\textwidth}
\begin{cardidea}{Contexto}{\faWater}
{\footnotesize En una autopista mojada por la lluvia, la \textbf{resistencia al avance} de un automóvil (fricción con la película de agua) depende de la velocidad. Ingenieros de transporte midieron la fuerza de fricción $F$ (N) en varias velocidades $v$ (m/s):}\\[0.8mm]
\centering
{\footnotesize\renewcommand{\arraystretch}{1.25}\setlength{\tabcolsep}{7pt}\arrayrulecolor{Granate}
\begin{tabular}{c|cccc}
$v$ (m/s) & 0.0 & 0.5 & 1.0 & 1.5\\\hline
$F$ (N) & 0.0 & 12.5 & 21.0 & 28.0\\
\end{tabular}}\\[0.8mm]
{\footnotesize El trabajo es $W=\displaystyle\int_0^{1.5}F(v)\,dv.$}
\end{cardidea}
\end{column}
\begin{column}{0.46\textwidth}
\begin{cardidea}{Se pide}{\faTasks}
{\footnotesize
\textbf{a)\ (2.0 P)}\ Use la regla de \textbf{Simpson $3/8$} para estimar el trabajo en Joules cuando la velocidad varía de $0$ a $1.5$ m/s.\\[1.4mm]
\textbf{b)\ (1.0 P)}\ Estime $W$ usando la \textbf{regla del trapecio} con los mismos datos.\\[1.4mm]
\textbf{c)\ (1.0 P)}\ Calcule la \textbf{diferencia} entre los resultados de Simpson $3/8$ y el trapecio. Comente cuál es más preciso y por qué.}
\end{cardidea}
\end{column}
\end{columns}
\end{frame}

% ---- PII-3 solución ----
\begin{frame}{Problema 3 — Simpson $3/8$, trapecio y comparación}
\begin{columns}[T]
\begin{column}{0.49\textwidth}
\begin{cardgear}{(a) Simpson $3/8$ ($h=0.5$)}{\faCalculator}
{\footnotesize $h=\dfrac{1.5-0}{3}=0.5$. Con 4 puntos ($n=3$):
\[
W=\tfrac{3h}{8}\big[f_0+3f_1+3f_2+f_3\big].
\]
\[
W=\tfrac{3(0.5)}{8}\big[0+3(12.5)+3(21.0)+28\big]
\]
\[
=0.1875\,[0+37.5+63+28]=0.1875(128.5).
\]
\[
\boxed{W_{3/8}=24.09375\ \text{J}}
\]}
\end{cardgear}
\end{column}
\begin{column}{0.49\textwidth}
\begin{cardgear}{(b) Trapecio ($h=0.5$)}{\faCalculator}
{\footnotesize
\[
W=\tfrac{h}{2}\big[f_0+2f_1+2f_2+f_3\big].
\]
\[
=\tfrac{0.5}{2}\big[0+2(12.5)+2(21.0)+28\big]
\]
\[
=0.25\,[0+25+42+28]=0.25(95).
\]
\[
\boxed{W_{\text{trap}}=23.75\ \text{J}}
\]}
\end{cardgear}
\vspace{1mm}
\begin{cardnota}[\faCheckCircle]
{\footnotesize \textbf{(c)}\ $|W_{3/8}-W_{\text{trap}}|=|24.09375-23.75|=0.34375$ J. Simpson $3/8$ es más preciso: ajusta con \emph{cúbicas} frente a los segmentos lineales del trapecio.}
\end{cardnota}
\end{column}
\end{columns}
\end{frame}

% ############################################################
%  PARTE II — PROBLEMA 4
% ############################################################
\portadaseccion{II}{Problema 4 — Polea con resorte torsional}{Forma vectorial, Lipschitz y método de Euler modificado (Heun)}{\faCogs}

% ---- PII-4 enunciado ----
\begin{frame}{Problema 4 — Enunciado (sistema mecánico)}
\begin{columns}[T]
\begin{column}{0.5\textwidth}
\begin{cardidea}{Modelo físico}{\faCogs}
{\footnotesize Una \textbf{polea sólida} (disco) de masa $m_1=4.0$ kg y radio $r=1.0$ m está conectada a un \textbf{resorte torsional} de constante $c_F=16\ \mathrm{N\cdot m/rad}$. Una cuerda enrollada sostiene una masa $m_2=2.0$ kg (el otro extremo fijo); $g=10\ \mathrm{m/s^2}$. Ecuación de movimiento:
\[
\Big(\tfrac{m_1}{2}+m_2\Big)\ddot x+\tfrac{c_F}{r^{2}}x=m_2 g,
\]
con $x(0)=0,\ \dot x(0)=0$. Sustituyendo:
\[
4\ddot x+16x=20\ \Rightarrow\ \ddot x+4x=5.
\]}
\end{cardidea}
\end{column}
\begin{column}{0.48\textwidth}
\begin{cardidea}{Se pide}{\faTasks}
{\scriptsize
\textbf{a)\ (1.0 P)}\ Escriba el sistema en forma vectorial $\dot{\mathbf Z}=F(t,\mathbf Z)$, $\mathbf Z=(x,\dot x)^{\!\top}$. Demuestre que $F$ cumple el \textbf{teorema de unicidad} (constante de Lipschitz con norma $\infty$ de $J_F$). ¿Por qué importa?\\[0.8mm]
\textbf{b)\ (1.0 P)}\ Aplique dos pasos del método de \textbf{Euler modificado (Heun)} vectorial con $h=0.5$ s para aproximar $\mathbf Z(0.5)$ y $\mathbf Z(1.0)$.\\[0.8mm]
\textbf{c)\ (1.0 P)}\ Si $x(t)=1.25(1-\cos 2t)$, calcule el \textbf{error absoluto} en $t=0.5$ y $t=1.0$. Justifique cómo $h$ afecta la precisión.\\[0.8mm]
\textbf{d)\ (1.0 P)}\ Si se usara \textbf{Euler explícito} en vez del modificado, ¿cómo cambiaría la precisión? Justifique por estabilidad y orden.}
\end{cardidea}
\end{column}
\end{columns}
\end{frame}

% ---- PII-4(a) ----
\begin{frame}{Problema 4(a) — Forma vectorial y Lipschitz}
\begin{columns}[T]
\begin{column}{0.5\textwidth}
\begin{cardgear}{Sistema de primer orden}{\faCogs}
{\footnotesize Con $z_1=x$, $z_2=\dot x$ y $\ddot x=5-4x$:
\[
\dot{\mathbf Z}=F(t,\mathbf Z)=\begin{pmatrix}z_2\\ 5-4z_1\end{pmatrix},\quad \mathbf Z(0)=\begin{pmatrix}0\\0\end{pmatrix}.
\]
Jacobiano de $F$:
\[
J_F=\frac{\partial F}{\partial \mathbf Z}=\begin{pmatrix}0 & 1\\ -4 & 0\end{pmatrix}.
\]}
\end{cardgear}
\end{column}
\begin{column}{0.48\textwidth}
\begin{cardgear}{Constante de Lipschitz}{\faRulerHorizontal}
{\footnotesize Sumas por fila de $|J_F|$: $0+1=1$ y $4+0=4$.
\[
\lVert J_F\rVert_\infty=\max(1,4)=4=L.
\]}
\end{cardgear}
\vspace{1mm}
\begin{cardnota}[\faCheckCircle]
{\footnotesize $F$ es \textbf{Lipschitz} ($L=4$ acotado en el dominio): por el \textbf{teorema de existencia y unicidad}, la solución es \emph{única} y la aproximación numérica es \textbf{estable} en $t$.}
\end{cardnota}
\end{column}
\end{columns}
\end{frame}

% ---- PII-4(b) ----
\begin{frame}{Problema 4(b) — Euler modificado (Heun), $h=0.5$}
\begin{columns}[T]
\begin{column}{0.5\textwidth}
\begin{cardgear}{Paso 1: $\mathbf Z(0)\to\mathbf Z(0.5)$}{\faArrowRight}
{\footnotesize \textbf{Predictor} (Euler), $F(\mathbf Z_0)=(0,\,5)$:
\[
\mathbf Z^{(1,e)}=\mathbf Z_0+hF(\mathbf Z_0)=\begin{pmatrix}0\\2.5\end{pmatrix}.
\]
\textbf{Corrector}, $F(\mathbf Z^{(1,e)})=(2.5,\,5)$:
\[
\mathbf Z_1=\mathbf Z_0+\tfrac{h}{2}\big[F(\mathbf Z_0)+F(\mathbf Z^{(1,e)})\big]=\begin{pmatrix}0.625\\ 2.5\end{pmatrix}.
\]}
\end{cardgear}
\end{column}
\begin{column}{0.48\textwidth}
\begin{cardgear}{Paso 2: $\mathbf Z(0.5)\to\mathbf Z(1.0)$}{\faArrowRight}
{\footnotesize \textbf{Predictor}, $F(\mathbf Z_1)=(2.5,\,2.5)$:
\[
\mathbf Z^{(2,e)}=\mathbf Z_1+hF(\mathbf Z_1)=\begin{pmatrix}1.875\\ 3.75\end{pmatrix}.
\]
\textbf{Corrector}, $F(\mathbf Z^{(2,e)})=(3.75,\,-2.5)$:
\[
\mathbf Z_2=\mathbf Z_1+\tfrac{h}{2}\big[F(\mathbf Z_1)+F(\mathbf Z^{(2,e)})\big]=\begin{pmatrix}2.1875\\ 2.5\end{pmatrix}.
\]}
\end{cardgear}
\vspace{1mm}
\begin{cardnota}[\faCheckCircle]
{\footnotesize $\mathbf Z(0.5)\approx(0.625,\,2.5)$ y $\mathbf Z(1.0)\approx(2.1875,\,2.5)$.}
\end{cardnota}
\end{column}
\end{columns}
\end{frame}

% ---- PII-4(c)(d) ----
\begin{frame}{Problema 4(c, d) — Error absoluto y comparación de métodos}
\begin{columns}[T]
\begin{column}{0.52\textwidth}
\begin{cardgear}{Error frente a la solución analítica}{\faTable}
{\footnotesize $x(t)=1.25(1-\cos 2t)$,\ $\dot x(t)=2.5\sin 2t$:}\\[1mm]
\centering
{\scriptsize\renewcommand{\arraystretch}{1.25}\setlength{\tabcolsep}{4.5pt}\arrayrulecolor{Granate}
\begin{tabular}{c|cc|c}
$t$ & aprox. & exacto & $|$error$|$\\\hline
\multicolumn{4}{l}{\itshape posición $x$}\\
0.5 & 0.6250 & 0.5746 & 0.0504\\
1.0 & 2.1875 & 1.7702 & 0.4173\\\hline
\multicolumn{4}{l}{\itshape velocidad $\dot x$}\\
0.5 & 2.5000 & 2.1037 & 0.3963\\
1.0 & 2.5000 & 2.2732 & 0.2268\\
\end{tabular}}
\end{cardgear}
\vspace{1mm}
\begin{cardnota}[\faExclamationTriangle]
{\scriptsize \textbf{Corrección.} La clave imprime $x(1.0)=2.416$; recalculando $x(1)=1.25(1-\cos 2)=1.7702$ (su propio error $0.417=|2.1875-1.7702|$ lo confirma). También $|$error$|$ en $t{=}0.5$ es $0.050$, no $0.104$.}
\end{cardnota}
\end{column}
\begin{column}{0.46\textwidth}
\begin{cardidea}{Efecto de $h$ y orden}{\faLightbulb}
{\footnotesize \textbf{(c)}\ Heun es de \textbf{orden 2}: error global $O(h^2)$. Con $h=0.5$ los errores son $\sim\!\tfrac14$; al reducir a $h=0.25$ caerían a $\approx0.06$ (factor $\sim4$).}
\end{cardidea}
\vspace{1mm}
\begin{cardidea}{(d) Euler explícito}{\faExclamationTriangle}
{\footnotesize Euler explícito es de \textbf{orden 1} ($O(h)$): \textbf{menos preciso} al mismo $h$. En sistemas \emph{oscilatorios} como este puede introducir \textbf{inestabilidades} si $h$ es grande.}
\end{cardidea}
\end{column}
\end{columns}
\end{frame}

% ============================================================
%  CIERRE
% ============================================================
\begin{frame}[plain]
\begin{tikzpicture}[remember picture, overlay]
  \fill[FondoOscuro] (current page.south west) rectangle (current page.north east);
  \node[color=IconoTenue] at ([xshift=-3.4cm,yshift=0.3cm]current page.east)
      {\fontsize{62}{62}\selectfont \faClipboardCheck};
  \fill[GranateVelo] (current page.north west)
      -- ([xshift=8.6cm]current page.north west)
      -- ([xshift=11.3cm]current page.south west)
      -- (current page.south west) -- cycle;
  \fill[GranateOscuro, opacity=0.75] (current page.north west)
      -- ([xshift=6.4cm]current page.north west)
      -- ([xshift=9.1cm]current page.south west)
      -- (current page.south west) -- cycle;
  \draw[Teal, line width=1.5pt] ([xshift=8.6cm]current page.north west)
      -- ([xshift=11.3cm]current page.south west);
  \draw[Naranja, line width=0.9pt] ([xshift=7.4cm]current page.north west)
      -- ([xshift=10.1cm]current page.south west);
  \node[anchor=west, text=white] at ([xshift=1.1cm,yshift=0.75cm]current page.west)
      {\fontsize{24}{28}\selectfont\bfseries ¡Examen resuelto!};
  \draw[white, line width=0.9pt] ([xshift=1.15cm,yshift=0.12cm]current page.west) -- ++(4.8cm,0);
  \node[anchor=west, text=white] at ([xshift=1.15cm,yshift=-0.55cm]current page.west)
      {{\color{white}\faAngleDoubleRight}\hspace{1.5mm}\small Métodos Numéricos (MB536) \textperiodcentered\ Examen Final 2025-1 \textperiodcentered\ Resuelto paso a paso};
\end{tikzpicture}
\end{frame}

\end{document}
